\documentclass[10pt, letterpaper]{article}

% Packages:
\usepackage[
    ignoreheadfoot, % set margins without considering header and footer
    top=2 cm, % seperation between body and page edge from the top
    bottom=2 cm, % seperation between body and page edge from the bottom
    left=2 cm, % seperation between body and page edge from the left
    right=2 cm, % seperation between body and page edge from the right
    footskip=1.0 cm, % seperation between body and footer
    % showframe % for debugging
]{geometry} % for adjusting page geometry
\usepackage[explicit]{titlesec} % for customizing section titles
\usepackage{tabularx} % for making tables with fixed width columns
\usepackage{array} % tabularx requires this
\usepackage[dvipsnames]{xcolor} % for coloring text
\definecolor{primaryColor}{RGB}{0, 79, 144} % define primary color
\usepackage{enumitem} % for customizing lists
\usepackage{fontawesome5} % for using icons
\usepackage{amsmath} % for math
\usepackage[
    pdftitle={Chunxu Han's CV},
    pdfauthor={Chunxu Han},
    pdfcreator={LaTeX with RenderCV},
    colorlinks=true,
    urlcolor=primaryColor
]{hyperref} % for links, metadata and bookmarks
\usepackage[pscoord]{eso-pic} % for floating text on the page
\usepackage{calc} % for calculating lengths
\usepackage{bookmark} % for bookmarks
\usepackage{lastpage} % for getting the total number of pages
\usepackage{changepage} % for one column entries (adjustwidth environment)
\usepackage{paracol} % for two and three column entries
\usepackage{ifthen} % for conditional statements
\usepackage{needspace} % for avoiding page brake right after the section title
\usepackage{iftex} % check if engine is pdflatex, xetex or luatex
\usepackage{multicol}
\usepackage{datetime2}

% Ensure that generate pdf is machine readable/ATS parsable:
\ifPDFTeX
    \input{glyphtounicode}
    \pdfgentounicode=1
    \usepackage[T1]{fontenc}
    \usepackage[utf8]{inputenc}
    \usepackage{lmodern}
\fi

\usepackage[default, type1]{sourcesanspro}

% Some settings:
\AtBeginEnvironment{adjustwidth}{\partopsep0pt} % remove space before adjustwidth environment
\pagestyle{empty} % no header or footer
\setcounter{secnumdepth}{0} % no section numbering
\setlength{\parindent}{0pt} % no indentation
\setlength{\topskip}{0pt} % no top skip
\setlength{\columnsep}{0.15cm} % set column seperation
\makeatletter
\let\ps@customFooterStyle\ps@plain % Copy the plain style to customFooterStyle
\patchcmd{\ps@customFooterStyle}{\thepage}{
    \color{gray}\textit{\small Chunxu Han - Page \thepage{} of \pageref*{LastPage}}
}{}{} % replace number by desired string
\makeatother
\pagestyle{customFooterStyle}

\titleformat{\section}{
    \needspace{4\baselineskip}
    \Large\color{primaryColor}
}{
}{
}{
    \textbf{#1}\hspace{0.15cm}\titlerule[0.8pt]\hspace{-0.1cm}
}[] % section title formatting

\titlespacing{\section}{
    -1pt
}{
    0.3 cm
}{
    0.2 cm
} % section title spacing

\newenvironment{highlights}{
    \begin{itemize}[
        topsep=0.10 cm,
        parsep=0.10 cm,
        partopsep=0pt,
        itemsep=0pt,
        leftmargin=0.4 cm + 10pt
    ]
}{
    \end{itemize}
} % new environment for highlights

\newenvironment{highlightsforbulletentries}{
    \begin{itemize}[
        topsep=0.10 cm,
        parsep=0.10 cm,
        partopsep=0pt,
        itemsep=0pt,
        leftmargin=10pt
    ]
}{
    \end{itemize}
} % new environment for highlights for bullet entries

\newenvironment{onecolentry}{
    \begin{adjustwidth}{
        0.2 cm + 0.00001 cm
    }{
        0.2 cm + 0.00001 cm
    }
}{
    \end{adjustwidth}
} % new environment for one column entries

\newenvironment{twocolentry}[2][]{
    \onecolentry
    \def\secondColumn{#2}
    \setcolumnwidth{\fill, 4.5 cm}
    \begin{paracol}{2}
}{
    \switchcolumn \raggedleft \secondColumn
    \end{paracol}
    \endonecolentry
} % new environment for two column entries

\newenvironment{threecolentry}[3][]{
    \onecolentry
    \def\thirdColumn{#3}
    \setcolumnwidth{1 cm, \fill, 4.5 cm}
    \begin{paracol}{3}
    {\raggedright #2} \switchcolumn
}{
    \switchcolumn \raggedleft \thirdColumn
    \end{paracol}
    \endonecolentry
} % new environment for three column entries

\newenvironment{header}{
    \setlength{\topsep}{0pt}\par\kern\topsep\centering\color{primaryColor}\linespread{1.5}
}{
    \par\kern\topsep
} % new environment for the header

% save the original href command in a new command:
\let\hrefWithoutArrow\href

% new command for external links:
\renewcommand{\href}[2]{\hrefWithoutArrow{#1}{\ifthenelse{\equal{#2}{}}{ }{#2 }\raisebox{.15ex}{\footnotesize \faExternalLink*}}}


\begin{document}
    \newcommand{\AND}{\unskip
        \cleaders\copy\ANDbox\hskip\wd\ANDbox
        \ignorespaces
    }
    \newsavebox\ANDbox
    \sbox\ANDbox{}

    \begin{header}
        \raggedright
        \fontsize{28 pt}{28 pt}
        \textbf{Chunxu Han}

        \vspace{0.3 cm}

        \normalsize
        \mbox{{\footnotesize\faMapMarker*}\hspace*{0.13cm}Peder Skrams Gade 16A, København K, Denmark}%
        \kern 0.25 cm%
        \AND%
        \kern 0.25 cm%
        \mbox{\hrefWithoutArrow{mailto:s220311@dtu.dk}{{\footnotesize\faEnvelope[regular]}\hspace*{0.13cm}s220311@dtu.dk}}%
        \kern 0.25 cm%
        \AND%
        \kern 0.25 cm%
        \mbox{\hrefWithoutArrow{tel:+45 52 73 06 85}{{\footnotesize\faPhone*}\hspace*{0.13cm}+45 52 73 06 85}}%
        \kern 0.25 cm%
        \AND%
        \kern 0.25 cm%
        \mbox{\hrefWithoutArrow{https://www.linkedin.com/in/chunxu-han/}{{\footnotesize\faLinkedinIn}\hspace*{0.13cm}chunxu-han}}%
        \kern 0.25 cm%
        \AND%
        \kern 0.25 cm%
        \mbox{\hrefWithoutArrow{https://github.com/chunxubioinfor}{{\footnotesize\faGithub}\hspace*{0.13cm}chunxubioinfor}}%
    \end{header}

    \vspace{0.4 cm - 0.3 cm}

    \section{Relevant Experience}

    \begin{onecolentry}

    \textbf{Machine Learning for High-Dimensional Omics Data} \\
    My master's thesis at Novo Nordisk focused on developing and benchmarking computational methods to evaluate batch-to-batch variability in single-cell RNA-seq data — a high-dimensional setting where model evaluation and generalizability across datasets are central challenges. I designed quantitative metrics from scratch, implemented a systematic benchmarking framework comparing multiple computational approaches, and explored machine learning methods to correlate transcriptomic variability with functional outcomes. This work resulted in \textbf{scBatch}, a Python package built around reproducibility and modular design.

    \vspace{0.3 cm}

    \textbf{Reproducible Pipelines and Research Software Development} \\
    Over the past two years, I contributed to the development and maintenance of \href{https://github.com/kvittingseerup/IsoformSwitchAnalyzeR}{IsoformSwitchAnalyzeR}, an R package for isoform-level analysis of long-read and single-cell sequencing data, which I helped bring to a published preprint. Across both roles, I built reproducible bioinformatics pipelines on HPC systems using Snakemake and Git, and worked extensively with large-scale omics datasets — including curation and re-quantification of 40,000+ proteomics samples. I am used to designing workflows that are not just functional, but transferable and well-documented.

    \vspace{0.3 cm}

    \textbf{Collaborative Research Across Academic and Industrial Settings} \\
    I have worked at the intersection of methods development and translational application — first in an academic isoform analysis group at DTU HealthTech, then in an industry-facing R\&D setting at Novo Nordisk, collaborating with cross-functional CMC and R\&D teams. I enjoy contributing to research that ultimately has a clinical or patient-relevant angle, and I find it motivating to work in environments where rigorous computational work connects to real biological questions.

    \end{onecolentry}

    \section{Education}

        \begin{twocolentry}{
    Lyngby, Denmark

    Sept 2023 -- March 2026
}
    \textbf{MSc in Bioinformatics and Systems Biology}, Technical University of Denmark
    \begin{highlights}
        \item Completed \textbf{Machine Learning} and \textbf{Deep Learning} courses alongside specialised bioinformatics training, building a strong foundation for developing and evaluating ML models in biological contexts.
    \end{highlights}
\end{twocolentry}

        \vspace{0.3 cm}

\begin{twocolentry}{
    Beijing, China

    Sept 2017 -- June 2021
}
    \textbf{BSc in Biotechnology}, Beijing Normal University
    \begin{highlights}
        \item Foundation in molecular biology, cell biology, and bioinformatics, with hands-on laboratory training.
    \end{highlights}
\end{twocolentry}

    \section{Work Experience}

        \begin{twocolentry}{
    Målev, Denmark

    Sep 2025 -- Mar 2026
}
    \textbf{Master Student}, Novo Nordisk

    \vspace{0.1cm}

    Thesis: developing computational frameworks to evaluate batch-to-batch variability in single-cell RNA-seq data for cell therapy quality control.

    \begin{highlights}
        \item \textbf{Benchmarking and Methodology Development:} Designed and implemented a systematic benchmarking framework comparing multiple computational tools for assessing batch comparability in cell-type composition and transcriptional profiles.

        \vspace{0.1cm}

        \item \textbf{Machine Learning:} Explored ML approaches to correlate transcriptomic variability with functional outcomes, contributing to predictive quality assessment in a GMP-adjacent setting.

        \vspace{0.1cm}

        \item \textbf{Software:} Developed \textbf{scBatch}, a Python package providing a structured and reproducible framework for batch variability assessment, with a focus on modular design and clear interfaces.

        \vspace{0.1cm}

        \item \textbf{Collaboration:} Worked with cross-functional R\&D and CMC teams to ensure biological interpretability and industrial relevance of results.
    \end{highlights}
\end{twocolentry}

        \vspace{0.3 cm}

        \begin{twocolentry}{
    Lyngby, Denmark

    Oct 2023 -- July 2025
}
    \textbf{Bioinformatics Research Assistant}, DTU HealthTech

    \vspace{0.1cm}

    Two years in the Isoform Analysis Group, contributing to both software development and large-scale data analysis.

    \begin{highlights}
        \item \textbf{Software Development:} Contributed to developing and updating \href{https://github.com/kvittingseerup/IsoformSwitchAnalyzeR}{IsoformSwitchAnalyzeR} v2 (R), supporting isoform-level analysis of long-read and single-cell RNA-seq data. The package is now available as a preprint on bioRxiv.

        \vspace{0.1cm}

        \item \textbf{Bioinformatics Pipelines:} Built and optimised reproducible pipelines on HPC infrastructure using Snakemake and Git; curated and re-quantified 40,000+ proteomics samples from GPMdb following FAIR data principles.

        \vspace{0.1cm}

        \item \textbf{Collaboration:} Supported PhD students, postdocs, and master's students with data preprocessing, pipeline troubleshooting, and result interpretation.
    \end{highlights}
\end{twocolentry}

        \vspace{0.3 cm}

        \begin{twocolentry}{
    Shanghai, China

    Feb 2023 -- May 2023
}
    \textbf{Consulting Intern}, Illuminera, IQVIA

    \begin{highlights}
        \item Conducted quantitative market and competitor analysis for pharmaceutical clients; delivered data-driven strategic reports using Power BI.
    \end{highlights}
\end{twocolentry}

        \vspace{0.3 cm}

        \begin{twocolentry}{
    Beijing, China

    Oct 2020 -- Nov 2020
}
    \textbf{QC Intern}, SINOVAC BIOTECH LTD.

    \begin{highlights}
        \item Worked in a GMP-regulated Quality Control lab; techniques include HP-LC, ELISA, SDS-PAGE, Lowry Protein Assay.
    \end{highlights}
\end{twocolentry}

\section*{Publications}

\begin{itemize}
  \item \textbf{Han, C.}, Gilis, J., Iriondo Delgado, E., Clement, L., \& Vitting-Seerup, K.\\
  \textit{IsoformSwitchAnalyzeR v2: Analysis of Functional Isoform Changes in Long-read and Single-cell Sequencing Data}.\\
  bioRxiv (2025). \href{https://doi.org/10.64898/2025.12.08.693027}{doi:10.64898/2025.12.08.693027}

  \item Mikkelsen, J. M., \textbf{Han, C.}, Kanakoglou, D. S., Tiberi, S., \& Vitting-Seerup, K.\\
  \textit{Beyond the Gene in Genetics: How Isoform-Resolved Analysis Empowers the Study of Both Common and Rare Genetic Variation}.\\
  medRxiv (2025). \href{https://doi.org/10.1101/2025.03.31.25324931}{doi:10.1101/2025.03.31.25324931}
\end{itemize}

    \section{Skills}

        \begin{onecolentry}
            \textbf{Programming:} Python, R, Bash, SQL, Snakemake, Docker
        \end{onecolentry}

        \vspace{0.2 cm}

        \begin{onecolentry}
            \textbf{Bioinformatics:} Single-cell RNA-seq, isoform analysis, long-read sequencing, HPC pipelines, batch effect methods
        \end{onecolentry}

        \vspace{0.2 cm}

        \begin{onecolentry}
            \textbf{Machine Learning:} Supervised/unsupervised learning, deep learning, model evaluation and benchmarking, foundation models for omics
        \end{onecolentry}

        \vspace{0.2 cm}

        \begin{onecolentry}
            \textbf{Languages:} English (proficient), Chinese (native), Danish (basic)
        \end{onecolentry}

        \vspace{0.2 cm}

        \begin{onecolentry}
            \textbf{Hobbies:}
            I've been playing \textbf{badminton} since an early age. I also enjoy \textbf{traveling} around Europe, \textbf{hiking}, and \textbf{sailing} — last summer we sailed to Ven and Hallands Väderö, which was a completely new kind of freedom.
        \end{onecolentry}

\end{document}
